\documentclass[11pt]{article}

\usepackage[margin=1in]{geometry}
\usepackage{amsmath,amssymb,amsthm}
\usepackage{enumitem}
\usepackage{hyperref}
\usepackage{xcolor}
\usepackage{booktabs}

\hypersetup{colorlinks=true, linkcolor=blue, urlcolor=blue, citecolor=blue}

\title{TOAE209/TOAH209 -- Design and Analysis of Algorithms\\[4pt]
\large Week 14: Theoretical Questions}
\author{Foreign Trade University}
\date{}

\begin{document}
\maketitle

\section*{Instructions}

Part A contains 18 questions requiring written explanations or short proofs about
approximation ratios and local search. Part B is a 10-question quiz (true/false and
multiple choice). Attempt every question on your own before consulting Part C (Solutions)
at the end of this document.

\section*{Part A: Theory Questions}

\begin{enumerate}[label=\textbf{A\arabic*.}, leftmargin=2.2em]

\item Define what it means for a polynomial-time algorithm $A$ to be a
\emph{$\rho$-approximation algorithm}, separately for minimization and maximization
problems. Why is $\rho = 1$ the best possible value, and what does a larger $\rho$ mean?

\item Every approximation proof this week follows the same pattern: instead of comparing
the algorithm's value $A(I)$ directly to $\mathrm{OPT}(I)$, we compare it to a
\emph{computable bound} $B(I)$. Explain why we cannot simply compare to
$\mathrm{OPT}(I)$ directly, and state the two chained inequalities the proof must
establish (for a minimization problem).

\item Distinguish a \emph{constant-factor} approximation, a \emph{PTAS}, and an
\emph{FPTAS}. Order them from weakest to strongest guarantee, and explain the difference
between a PTAS and an FPTAS in terms of how the running time depends on $\varepsilon$.

\item State the Minimum Vertex Cover problem. Then describe the maximal-matching algorithm
and prove that the set it returns is in fact a vertex cover.

\item For the maximal-matching Vertex Cover algorithm, prove the lower bound $\mathrm{OPT}
\ge |M|$, where $M$ is the matching the algorithm builds. Why is it essential that the
edges of $M$ are pairwise vertex-disjoint?

\item Combine the previous two facts into the full $2$-approximation proof for Vertex
Cover. Explain in one sentence why adding \emph{both} endpoints of each matched edge
(rather than just one) is what makes the analysis work.

\item State the Minimum Set Cover problem and the greedy algorithm. Then prove that an
element $e$ first covered when greedy selects a set covering $t$ uncovered elements is
``charged'' exactly $1/t$, and that the greedy cover size equals $\sum_{e} \mathrm{cost}(e)$.

\item Complete the analysis of greedy Set Cover: prove that the elements of any single
optimal set $S_j$ receive total charge at most $H_{|S_j|} \le H_n$, and conclude that
greedy is an $H_n = O(\log n)$ approximation. (State clearly where you use the greedy
choice rule.)

\item State the Load Balancing (makespan) problem. Prove the two lower bounds on the
optimal makespan: $\mathrm{OPT} \ge \frac{1}{m}\sum_j t_j$ and $\mathrm{OPT} \ge \max_j
t_j$.

\item Prove that greedy list scheduling is a $2$-approximation for load balancing. Identify
exactly where each of the two lower bounds from Q9 is used.

\item State (without full proof) the improved guarantee for Longest-Processing-Time (LPT)
scheduling. Explain intuitively why sorting jobs in \emph{decreasing} order helps, and why
the ``last/critical job'' on the bottleneck machine is small when there are at least two
jobs on it.

\item State the $k$-center problem and the greedy farthest-point algorithm. Why is the
\emph{metric} (triangle-inequality) assumption necessary -- what goes wrong without it?

\item Prove that greedy farthest-point selection is a $2$-approximation for $k$-center.
Your proof should produce $k+1$ sites that are pairwise far apart and apply pigeonhole over
the $k$ optimal centers.

\item State the Metric TSP problem and the MST-based algorithm. Prove the lower bound
$w(T) \le \mathrm{OPT}$ relating the MST weight to the optimal tour.

\item Complete the MST-based Metric TSP proof: explain the doubling/Eulerian-circuit and
shortcutting steps, and show why shortcutting cannot increase the tour length. State the
final $2$-approximation bound. How does Christofides improve the constant?

\item State the 0/1 Knapsack problem and the FPTAS via rounding-and-scaling. With scaling
factor $\mu = \varepsilon\, p_{\max}/n$, prove the quality guarantee $\sum_{i\in S} p_i \ge
(1-\varepsilon)\mathrm{OPT}$, and explain why the scaled DP runs in time polynomial in
$n$ and $1/\varepsilon$.

\item Describe the local-search framework: feasible solution, neighbourhood, improving
move, local optimum. For Max-Cut with the single-vertex-flip neighbourhood, prove that
every local optimum cuts at least $m/2$ edges, and conclude it is a $2$-approximation.

\item Hill climbing can get stuck at a local optimum that is far from global. Explain how
\emph{random restarts} and \emph{simulated annealing} each address this. For simulated
annealing, state the acceptance probability for a worsening move and describe the role of
the temperature schedule.

\end{enumerate}

\newpage
\section*{Part B: Quiz}

\begin{enumerate}[label=\textbf{B\arabic*.}, leftmargin=2.2em]

\item \textbf{(Multiple choice)} For a \emph{minimization} problem, a $2$-approximation
algorithm guarantees that its output value is:
\begin{enumerate}[label=(\alph*)]
    \item at most $\tfrac12$ of the optimum
    \item at most $2$ times the optimum
    \item at least $2$ times the optimum
    \item exactly the optimum
\end{enumerate}

\item \textbf{(True/False)} The Vertex Cover algorithm that takes both endpoints of a
maximal matching is a $2$-approximation, even though it never computes the optimal cover.

\item \textbf{(Short answer)} In the maximal-matching Vertex Cover algorithm, if the
matching $M$ has $7$ edges, what is the size of the cover returned, and what is the best
lower bound this gives on $\mathrm{OPT}$?

\item \textbf{(Multiple choice)} The greedy Set Cover algorithm achieves an approximation
ratio of:
\begin{enumerate}[label=(\alph*)]
    \item $2$
    \item $\tfrac32$
    \item $H_n = O(\log n)$
    \item $1 + \varepsilon$
\end{enumerate}

\item \textbf{(True/False)} Greedy list scheduling for load balancing assigns each job to
the machine of currently \emph{smallest} load.

\item \textbf{(Short answer)} For $m = 3$ machines and total processing time $\sum_j t_j =
30$ with largest job $\max_j t_j = 7$, what is the best lower bound on the optimal makespan
you can state from the two standard bounds?

\item \textbf{(Multiple choice)} The greedy farthest-point $k$-center algorithm relies on
which property of the distance function for its $2$-approximation guarantee?
\begin{enumerate}[label=(\alph*)]
    \item symmetry only
    \item the triangle inequality
    \item that all distances are integers
    \item nothing -- it works for arbitrary distances
\end{enumerate}

\item \textbf{(True/False)} For Metric TSP, the weight of a minimum spanning tree is a
lower bound on the length of the optimal tour.

\item \textbf{(Short answer)} An FPTAS for Knapsack returns a solution of value at least
$(1-\varepsilon)\mathrm{OPT}$. If $\varepsilon = 0.1$ and $\mathrm{OPT} = 200$, what is the
worst-case value the FPTAS may return?

\item \textbf{(True/False)} At a single-vertex-flip local optimum of Max-Cut, every vertex
has at least half of its incident edges crossing the cut, which forces the cut to contain
at least $m/2$ edges.

\end{enumerate}

\newpage
\section*{Part C: Solutions}

\subsection*{Part A}

\begin{enumerate}[label=\textbf{A\arabic*.}, leftmargin=2.2em]

\item $A$ is a \textbf{$\rho$-approximation} (with $\rho \ge 1$) if for every instance
$I$: for \emph{minimization}, $A(I) \le \rho\cdot\mathrm{OPT}(I)$; for \emph{maximization},
$A(I) \ge \tfrac{1}{\rho}\mathrm{OPT}(I)$. The value $\rho = 1$ means the algorithm is
exactly optimal on every input (the best possible, since by definition $A(I)$ can be no
better than $\mathrm{OPT}(I)$). A larger $\rho$ means a weaker guarantee -- the solution
may be up to $\rho$ times worse (minimization) or only a $1/\rho$ fraction as good
(maximization) as optimal.

\item We cannot compare to $\mathrm{OPT}(I)$ directly because computing $\mathrm{OPT}(I)$
is the NP-hard problem we are trying to avoid -- there is no efficient way to know its
value. Instead we find a quantity $B(I)$ we \emph{can} reason about and that bounds the
optimum, then prove two chained inequalities: (1) $B(I) \le \mathrm{OPT}(I)$ (the bound is
valid) and (2) $A(I) \le \rho\cdot B(I)$ (the algorithm is close to the bound). Chaining
gives $A(I) \le \rho\cdot B(I) \le \rho\cdot\mathrm{OPT}(I)$, all without computing
$\mathrm{OPT}(I)$.

\item A \textbf{constant-factor} approximation guarantees $\rho = c$ for a fixed constant
independent of the input (e.g.\ $\rho = 2$). A \textbf{PTAS} (polynomial-time
approximation scheme) achieves $\rho = 1 + \varepsilon$ for any fixed $\varepsilon > 0$, in
time polynomial in $n$ but possibly exponential in $1/\varepsilon$ (e.g.\
$n^{O(1/\varepsilon)}$). An \textbf{FPTAS} achieves $1 + \varepsilon$ in time polynomial in
\emph{both} $n$ \emph{and} $1/\varepsilon$ (e.g.\ $O(n^3/\varepsilon)$). Weakest to
strongest: constant-factor $<$ PTAS $<$ FPTAS. The distinction between PTAS and FPTAS is
precisely the dependence on $\varepsilon$: an FPTAS stays efficient even as $\varepsilon
\to 0$, while a PTAS may blow up.

\item \textbf{Minimum Vertex Cover:} given $G = (V, E)$, find a smallest $S \subseteq V$
such that every edge has at least one endpoint in $S$. \textbf{Algorithm:} scan the edges;
whenever an edge $(u,v)$ has neither endpoint yet in $S$, add \emph{both} $u$ and $v$ to
$S$ (and record $(u,v)$ in a matching $M$). \textbf{$S$ is a cover:} consider any edge
$(x,y)$. If it were uncovered ($x,y \notin S$), then when the scan reached $(x,y)$ neither
endpoint was in $S$, so the algorithm would have added them -- contradiction. Hence every
edge has an endpoint in $S$.

\item The edges of $M$ are pairwise vertex-disjoint: once $(u,v)$ is added, both $u$ and
$v$ are in $S$, so any later edge sharing $u$ or $v$ fails the ``neither endpoint in $S$''
test and is not added to $M$. Now \emph{any} vertex cover $C$ must cover each edge of $M$,
i.e.\ contain at least one endpoint of each $M$-edge. Because the $M$-edges share no
endpoints, these chosen endpoints are all \emph{distinct}, so $|C| \ge |M|$. Applying this
to an optimal cover gives $\mathrm{OPT} \ge |M|$. Disjointness is essential: if two
$M$-edges shared an endpoint, a single vertex could cover both, and we could not conclude
$|C| \ge |M|$.

\item The algorithm returns $S$ = endpoints of $M$, so $|S| = 2|M|$ (each edge contributes
two distinct vertices). By Q5, $\mathrm{OPT} \ge |M|$. Therefore $|S| = 2|M| \le
2\,\mathrm{OPT}$, a $2$-approximation. Adding both endpoints (rather than one) is what
guarantees $S$ is a cover \emph{and} that $|S| = 2|M|$ ties cleanly to the lower bound
$|M|$; choosing one endpoint per matched edge would not guarantee coverage of all edges.

\item \textbf{Minimum Set Cover:} given universe $U$ ($|U| = n$) and subsets $S_1, \ldots,
S_m$ with $\bigcup S_i = U$, find the fewest sets whose union is $U$. \textbf{Greedy:}
while uncovered elements remain, pick the set covering the most currently-uncovered
elements; add it; repeat. \textbf{Charging:} when greedy picks a set covering $t$
previously-uncovered elements, we assign cost $1/t$ to each of those $t$ elements -- the
single set used is ``paid for'' by its newly-covered elements, total $t \cdot (1/t) = 1$.
Each element is charged exactly once, at the moment it is first covered. Summing,
$\sum_{e} \mathrm{cost}(e)$ telescopes to the total number of sets greedy picks.

\item Fix an optimal set $S_j$ and order its elements $e_1, \ldots, e_{|S_j|}$ by the time
greedy first covers them. Just before greedy covers $e_k$, the elements $e_k, \ldots,
e_{|S_j|}$ are all still uncovered, so $S_j$ \emph{itself} could cover at least $|S_j| - k
+ 1$ uncovered elements. \emph{Here we use the greedy rule:} greedy picks the set covering
the \emph{most} uncovered elements, so the set it actually picks covers at least $|S_j| - k
+ 1$, giving each newly-covered element (including $e_k$) cost $\le \frac{1}{|S_j|-k+1}$.
Summing, $\sum_{e\in S_j}\mathrm{cost}(e) \le \sum_{k=1}^{|S_j|}\frac{1}{|S_j|-k+1} =
H_{|S_j|} \le H_n$. Finally, let $\mathcal O$ be an optimal cover; every element lies in
some $S_j \in \mathcal O$, so $|\text{greedy}| = \sum_e \mathrm{cost}(e) \le \sum_{S_j \in
\mathcal O}\sum_{e \in S_j}\mathrm{cost}(e) \le |\mathcal O|\cdot H_n = \mathrm{OPT}\cdot
H_n$.

\item \textbf{Load Balancing:} $m$ identical machines, jobs with times $t_1, \ldots, t_n$;
assign each job to a machine; load $L_k$ = sum of times on machine $k$; minimize makespan
$\max_k L_k$. \textbf{Bound 1 (average):} the total work $\sum_j t_j$ is distributed over
$m$ machines, so by averaging some machine has load $\ge \frac1m\sum_j t_j$, hence the
makespan is $\ge \frac1m\sum_j t_j$; this holds for every assignment, including the
optimal one, so $\mathrm{OPT} \ge \frac1m\sum_j t_j$. \textbf{Bound 2 (max job):} the
machine running the largest job has load $\ge \max_j t_j$, so $\mathrm{OPT} \ge \max_j
t_j$.

\item Let machine $k$ achieve the makespan $L$, and let $j$ be the last job placed on it.
Just before placing $j$, machine $k$ had load $L - t_j$, and since greedy always places on
a least-loaded machine, \emph{every} machine then had load $\ge L - t_j$. Summing over
$m$ machines, the work placed so far is $\ge m(L - t_j)$, so $L - t_j \le \frac1m\sum_i t_i
\le \mathrm{OPT}$ (\emph{using Bound 1}). Also $t_j \le \max_i t_i \le \mathrm{OPT}$
(\emph{using Bound 2}). Adding, $L = (L - t_j) + t_j \le 2\,\mathrm{OPT}$.

\item LPT (sort jobs in decreasing order, then greedy) achieves makespan $\le \frac43
\mathrm{OPT}$ (more precisely $(\frac43 - \frac{1}{3m})\mathrm{OPT}$). Intuition: large
jobs are the hard part; scheduling them first, while many small jobs remain to ``fill the
gaps,'' lets the small jobs balance the machines afterward. If the critical (last) job on
the bottleneck machine is \emph{not alone} on that machine, then -- because jobs are placed
in decreasing order and at least two jobs are on the (most-loaded) machine -- there are at
least $m+1$ jobs of size $\ge t_j$, which forces $t_j \le \frac13\mathrm{OPT}$ (otherwise
the total work exceeds what $m$ machines can hold within $\mathrm{OPT}$). Then $L = (L -
t_j) + t_j \le \mathrm{OPT} + \frac13\mathrm{OPT} = \frac43\mathrm{OPT}$. The remaining
case (critical job alone) is handled directly and turns out optimal.

\item \textbf{$k$-center:} given sites $P$ in a metric space and integer $k$, choose $k$
centers $C$ minimizing $r(C) = \max_{p}\min_{c\in C} d(p,c)$. \textbf{Greedy:} start with
any center; repeatedly add the site \emph{farthest} from the current center set, until $k$
centers are chosen. The triangle inequality is necessary because the proof bounds the
distance between two sites served by the same optimal center via $d(a,b) \le d(a,c^\ast) +
d(c^\ast,b)$. Without it, two sites both close to a common center could still be
arbitrarily far from each other, so ``farthest point'' selection carries no guarantee, and
indeed no constant-factor approximation is possible for general (non-metric) distances.

\item Let $r = r(C)$ be the radius returned, realized by some site $p$ at distance $r$ from
its nearest center. Consider the $k$ chosen centers together with $p$: that is $k+1$ sites.
By the greedy rule, each center was chosen as the point farthest from the previously chosen
centers, and the leftover distance of $p$ is $r$; consequently every pair among these
$k+1$ sites is at distance $\ge r$. Suppose for contradiction $r > 2r^\ast$. There are
only $k$ optimal centers, so by pigeonhole two of the $k+1$ sites, say $a$ and $b$, share
the same optimal center $c^\ast$, giving $d(a,c^\ast), d(b,c^\ast) \le r^\ast$. By the
triangle inequality $d(a,b) \le 2r^\ast < r$, contradicting $d(a,b) \ge r$. Hence $r \le
2r^\ast$: a $2$-approximation.

\item \textbf{Metric TSP:} $n$ cities with symmetric distances satisfying the triangle
inequality; find a minimum-length tour visiting each city once and returning to start.
\textbf{Algorithm:} build an MST $T$, then output the preorder (DFS) traversal of $T$ as
the visiting order. \textbf{Lower bound:} take an optimal tour and delete any one edge.
What remains is a Hamiltonian \emph{path}, which is a spanning tree; its weight is $\le$
the tour length. Since the MST is the lightest spanning tree, $w(T) \le \text{(this path)}
\le \mathrm{OPT}$.

\item Double every edge of $T$ to get $2T$; now all degrees are even, so $2T$ has an
Eulerian circuit of total length exactly $2w(T)$. Traverse it and \emph{shortcut}: when
about to revisit a city, jump directly to the next unvisited one -- this yields the
preorder order and a valid tour. Each shortcut replaces a path between two cities by the
direct edge, which by the triangle inequality is no longer; so the tour length $\le 2w(T)$.
Combined with $w(T) \le \mathrm{OPT}$, the tour $\le 2w(T) \le 2\mathrm{OPT}$, a
$2$-approximation. \textbf{Christofides} improves the constant to $\frac32$ by adding a
minimum-weight perfect matching on the \emph{odd-degree} vertices of $T$ (instead of
doubling all edges), producing an Eulerian graph of weight $\le \mathrm{OPT} +
\frac12\mathrm{OPT}$.

\item \textbf{0/1 Knapsack:} items with profits $p_i$, weights $w_i$, capacity $W$; pick a
subset of weight $\le W$ maximizing total profit. \textbf{FPTAS:} set $\mu =
\varepsilon\,p_{\max}/n$, round profits down to $\hat p_i = \lfloor p_i/\mu\rfloor$, and
solve exactly under $\hat p_i$ by the profit-indexed DP. \textbf{Quality:} rounding gives
$\mu\hat p_i \le p_i \le \mu\hat p_i + \mu$. Let $O$ be optimal under true profits and $S$
the DP's set (optimal under $\hat p$), so $\sum_{i\in S}\hat p_i \ge \sum_{i\in O}\hat
p_i$. Then $\sum_{i\in S}p_i \ge \mu\sum_{i\in S}\hat p_i \ge \mu\sum_{i\in O}\hat p_i \ge
\sum_{i\in O}p_i - n\mu = \mathrm{OPT} - \varepsilon p_{\max} \ge (1-\varepsilon)
\mathrm{OPT}$, using $\mathrm{OPT}\ge p_{\max}$. \textbf{Running time:} each $\hat p_i \le
p_{\max}/\mu = n/\varepsilon$, so total rounded profit $\le n^2/\varepsilon$, and the DP
runs in $O(n\cdot\sum_i\hat p_i) = O(n^3/\varepsilon)$ -- polynomial in $n$ and
$1/\varepsilon$.

\item \textbf{Framework:} start from any feasible solution; the \emph{neighbourhood} is
the set of solutions reachable by a small ``move''; repeatedly move to a strictly better
neighbour (\emph{improving move}); stop at a \emph{local optimum}, where no neighbour is
better. \textbf{Max-Cut (flip neighbourhood):} for each vertex $v$, let $d_{\text{cut}}(v)$
and $d_{\text{same}}(v)$ be its numbers of cut / same-side edges (summing to $\deg v$).
Flipping $v$ changes the cut by $d_{\text{same}}(v) - d_{\text{cut}}(v)$. At a local
optimum flipping never helps, so $d_{\text{same}}(v) \le d_{\text{cut}}(v)$, i.e.\
$d_{\text{cut}}(v) \ge \frac12\deg v$ for all $v$. Summing, $2\cdot(\text{cut}) = \sum_v
d_{\text{cut}}(v) \ge \frac12\sum_v \deg v = \frac12(2m) = m$, so the cut $\ge m/2$. Since
$\mathrm{OPT} \le m$, this is a $2$-approximation.

\item \textbf{Random restarts:} run hill climbing from many random initial solutions and
keep the best local optimum found; since different starts lie in different ``basins of
attraction,'' more restarts raise the probability of reaching the global optimum.
\textbf{Simulated annealing:} occasionally accept a \emph{worsening} move to escape local
optima, with acceptance probability $e^{-\Delta/T}$ for a move that worsens the objective
by $\Delta > 0$, where $T$ is a temperature. The schedule starts $T$ high (most moves
accepted, broad exploration) and gradually lowers it toward $0$ (only improving moves
accepted, settling into a good solution); a sufficiently slow schedule converges to the
global optimum in the limit.

\end{enumerate}

\subsection*{Part B}

\begin{enumerate}[label=\textbf{B\arabic*.}, leftmargin=2.2em]

\item \textbf{(b) at most $2$ times the optimum.} For minimization, a $\rho$-approximation
guarantees $A(I) \le \rho\cdot\mathrm{OPT}(I)$.

\item \textbf{True.} It returns the endpoints of a maximal matching $M$, of size $2|M| \le
2\,\mathrm{OPT}$, never computing the optimal cover -- yet the matching is a valid lower
bound.

\item Cover size $= 2 \times 7 = \mathbf{14}$ vertices; lower bound $\mathrm{OPT} \ge |M|
= \mathbf{7}$.

\item \textbf{(c) $H_n = O(\log n)$.} This is the greedy Set Cover guarantee, and (up to
constants) the best possible unless $\mathrm{P}=\mathrm{NP}$.

\item \textbf{True.} Greedy list scheduling assigns each job in turn to a machine of
currently minimum load.

\item Bound 1: $\frac1m\sum_j t_j = 30/3 = 10$. Bound 2: $\max_j t_j = 7$. The best
(largest) lower bound is $\max(10, 7) = \mathbf{10}$.

\item \textbf{(b) the triangle inequality.} The $2$-approximation proof uses $d(a,b) \le
d(a,c^\ast) + d(c^\ast,b)$.

\item \textbf{True.} Deleting an edge from any tour leaves a spanning tree of weight $\le$
the tour, so the minimum spanning tree weight $\le \mathrm{OPT}$.

\item $(1 - 0.1)\times 200 = 0.9 \times 200 = \mathbf{180}$.

\item \textbf{True.} This is exactly the local-optimum condition $d_{\text{cut}}(v) \ge
\frac12\deg v$ for every vertex, which sums to a cut of size $\ge m/2$.

\end{enumerate}

\end{document}
